\chapter{Présentation de l'organisme d'accueil}
\label{chap:organisme}

\section{Le Groupe OCP}

Le Groupe OCP (Office Chérifien des Phosphates) est le premier exportateur mondial de
phosphates et de leurs dérivés. Son activité repose sur une ressource dont le Maroc détient
l'essentiel des réserves mondiales : la roche phosphatée.

L'activité du groupe couvre l'ensemble de la chaîne de valeur, de l'extraction minière à la
commercialisation d'engrais :

\begin{itemize}
  \item \textbf{extraction} de la roche phosphatée sur les sites miniers ;
  \item \textbf{valorisation chimique}, par transformation de la roche en acide phosphorique
        puis en engrais ;
  \item \textbf{commercialisation} sur les marchés internationaux.
\end{itemize}

Cette intégration verticale distingue le groupe de la plupart de ses concurrents et
constitue, on le verra, la source même du problème traité dans ce rapport.

\section{La plateforme chimique de Jorf Lasfar}

Située sur la côte atlantique près d'El Jadida, Jorf Lasfar est la plus grande plateforme
chimique du groupe. Elle rassemble sur un même site un ensemble d'unités interdépendantes.

\begin{table}[H]
\centering
\caption{Principales unités de la plateforme de Jorf Lasfar}
\label{tab:unites-jorf}
\begin{tabular}{@{}ll@{}}
\toprule
\textbf{Unité} & \textbf{Fonction} \\
\midrule
Ateliers sulfuriques & Production d'acide sulfurique \ce{H2SO4} \\
Ateliers phosphoriques & Production d'acide phosphorique \ce{H3PO4} --- \emph{périmètre du stage} \\
Ateliers d'engrais & Production de DAP, MAP, TSP, NPK \\
Coentreprises & EMAPHOS, IMACID, PMP, JFC \\
Utilités & Centrale thermique, station de dessalement, port minéralier \\
\bottomrule
\end{tabular}
\end{table}

\subsection{Une intégration qui crée l'interdépendance}

Sur un site intégré, les ateliers ne fonctionnent pas indépendamment : \emph{la sortie de
l'un est l'entrée de l'autre}. Un déséquilibre local se propage immédiatement à l'ensemble
de la chaîne. Une cuve pleine arrête la production amont, faute d'exutoire ; une cuve vide
arrête la production aval, faute d'alimentation.

\begin{encadre}
\textbf{L'origine du problème.} Coordonner quotidiennement ces flux --- décider qui produit
quoi, qui alimente qui, et dans quelles quantités --- est précisément l'objet de ce stage.
Le problème n'existerait pas sur un site où chaque atelier serait autonome.
\end{encadre}

\section{Les partenaires industriels du site}

Plusieurs des consommateurs modélisés dans ce travail sont des coentreprises, ce qui donne à
leurs demandes un caractère contractuel qu'il convient de garder à l'esprit.

\subsection{EMAPHOS --- Euro-Maroc Phosphore}

Coentreprise entre l'OCP, le belge Prayon et l'allemand Chemische Fabrik Budenheim, EMAPHOS
produit de l'\textbf{acide phosphorique purifié} de qualité alimentaire et technique, par
extraction liquide-liquide au solvant. Sa capacité, initialement de
\SI{140000}{\tonne}~\PdeuxOcinq{} par an, a été doublée à \SI{280000}{\tonne}~\PdeuxOcinq{}
par an à partir du quatrième trimestre 2022~\cite{ocp-emaphos}, soit de l'ordre de
\SIrange{380}{770}{\tonne}~\PdeuxOcinq{} par jour.

\begin{encadre}[title={Pourquoi cette capacité nous intéresse}]
La demande d'EMAPHOS dans le scénario de référence est de \SI{700}{\tonne} par jour. Ce
chiffre ne prend un sens réaliste que si l'unité de compte est la tonne de \PdeuxOcinq{}.
C'est l'une des trois preuves qui fondent la décision de modélisation la plus structurante
de ce travail (chapitre~\ref{chap:modelisation}, section~\ref{sec:unite}).
\end{encadre}

L'extraction au solvant sépare l'acide en deux phases : une phase organique chargée d'acide
purifié, et un \textbf{raffinat} aqueux concentrant les impuretés. Ce raffinat contient
encore une quantité importante de \PdeuxOcinq{} ; le rejeter constituerait une perte
économique majeure. Il est donc renvoyé en amont du circuit principal, où sa qualité
médiocre se dilue dans la masse. Ce retour représente \SI{40}{\percent} de la quantité
livrée et constitue, comme on le verra, le seul flux du système qui remonte d'un niveau de
concentration.

\subsection{Les autres consommateurs}

\begin{table}[H]
\centering
\caption{Consommateurs d'acide de la plateforme}
\label{tab:consommateurs}
\begin{tabular}{@{}lll@{}}
\toprule
\textbf{Consommateur} & \textbf{Nature} & \textbf{Acide consommé} \\
\midrule
U16, U116A, U116BC & Ateliers d'engrais internes & selon les recettes \\
EMAPHOS & Coentreprise OCP / Prayon / Budenheim & acide 54 NCL \\
IMACID & Coentreprise indo-marocaine & acide 29 standard \\
JFC 1--5 & \emph{Jorf Fertilizer Company} & acide 29 standard \\
MAPS & Unité partenaire & acide 29 standard \\
U53 & Unité interne & acide 54 clarifié \\
107DEF & Unité interne polyvalente & plusieurs qualités \\
\bottomrule
\end{tabular}
\end{table}

Du point de vue du modèle, chaque consommateur se caractérise par deux informations
seulement : les \textbf{types d'acide} qu'il accepte, et les \textbf{lignes} auxquelles il
est physiquement raccordé.

\section{Cadre du stage}

Ce stage s'inscrit dans le cursus d'ingénieur de l'INSEA, filière \filiere. Il mobilise
conjointement deux compétences de la formation :

\begin{itemize}
  \item la \textbf{recherche opérationnelle}, pour la modélisation et la résolution du
        problème d'optimisation ;
  \item le \textbf{génie logiciel}, pour produire un outil maintenable, testé et
        réutilisable plutôt qu'un script jetable.
\end{itemize}

Le livrable attendu comprend un modèle mathématique documenté, une implémentation
logicielle validée, des rapports d'exploitation au format imposé par le service, et le
présent rapport.
